\section{Exact scope, attribution and remaining source assertions}
The original constructions, their motivating representation problem and
their proposed further investigations belong to Leemans, Stokes and
Tranchida. The finite normal-pair existence result used here belongs to
D\'iaz Ramos, Molinier and Viruel and is proved in the complete first
application. No source correction is asserted as a new historical
result or as a newly solved question.

The finite examples and free-basis proofs here rework the explicit
phase197 manuscripts. The shadow theorem and its four-element projective
instance reproduce the phase196 argument with its separately recorded
independent review. They are not claimed as mathematics first proved
in this edition. The unrestricted panel-permutation construction additionally recovers the
weak free-group representation conclusions using a different compatible
inner action. This is not inferred from the finite signed subgroup, and
no historical priority is claimed.
The common recognition theorem was already used in
these earlier arguments; the present edition replaces its archived
reference by the complete accepted common paper and retains the full
F010 application. There is no unpaid premise in an old proof bank.

The classical automorphism facts for $S_n$ away from degree six and for
$A_5$ are explicit retained inputs. The positive finite projective repair
also retains the fundamental theorem for the full subspace geometry and
the standard automorphism-group identification for $\operatorname{PGL}(n,q)$,
$n\ge3$. The new ordered attachment and its full rigidity proof are supplied
here; these classical facts do not prove the attachment properties. The exceptional degree-six argument
is given here directly. The cited Atlas constructions for $A_6,A_7$ are
retained historical examples, not reconstructions of their individual
incidence tables. Their existence also follows from the complete finite
normal-pair theorem, without identifying its objects with those Atlas
models. The source's proposed orthogonal and Suzuki variants are not
proved constructions there and remain prospective investigations here.

Source Proposition 7.1 is the elementary coset incidence construction:
right multiplication preserves intersection, and intersecting cosets
have a common representative, proving transitivity on vertices and
incident pairs. The product-intersection and residual-generation criteria
quoted as Theorems 7.2 and 7.4 are retained classical results. For every
use in the free-basis family, the common paper gives direct tree and
residue proofs instead. The interpretation of ``independent'' in
Theorem 7.3 is a freely independent family, as required by its reduced-word
argument; the specific conjugate families satisfy this by the explicit
immersed finite core. No weaker irredundancy assumption is substituted.

The general Artin--Tits statement is a pointer to forthcoming work in
the source and is not a premise of these free-basis proofs. The background
Outer Space realization and maximum finite-subgroup discussion are not
needed to construct the signed permutation subgroup. In particular,
torsion-freeness of the inner subgroup by itself does not say that every
finite outer subgroup lifts to an isomorphic automorphism subgroup.
The finite signed actions proved here make no such inference. The source's
infinite-type and more restrictive finite-geometry questions remain
questions. The common connected-chamber splitting proposition supplies a valid
sufficient hypothesis. The complete thin-doubling construction and finite
certificate refute the unqualified splitting observation; they do not
resolve the remaining representation questions.

The inherited finite programs calculate all groups and flag sets in their
stated bounded models; the free-word program checks 45,846 reduced words.
The uniform proofs above do not infer infinite statements from those
samples. The projective program independently constructs all 2562 objects
and their fibers, but the written counting argument proves the displayed
obstruction without relying on the program's output. The full files are
retained as optional validation. The three basis-permutation Lean
certificates prove the action on free generators, transport of each
complementary generated subgroup and faithfulness. The seven inherited
shadow declarations prove the actual pointwise incidence kernel and its
triviality criterion; they do not prove the image or splitting of
restriction. Neither file formalizes the projective instance, free-basis
geometry, residual connectedness or normalizer. The twelve previously
accepted F010 certificates remain scoped as in the first application.

\begin{thebibliography}{9}
\bibitem{LST} D. Leemans, K. Stokes and P. Tranchida,
\emph{Groups represented by incidence geometries}, arXiv:2506.10464v1,
12 June 2025.
\bibitem{DMV} A. D\'iaz Ramos, R. Molinier and A. Viruel,
\emph{Every finite group is represented by a finite incidence geometry},
arXiv:2511.02410v2, 16 December 2025.
\bibitem{BC} F. Buekenhout and A. M. Cohen,
\emph{Diagram Geometry}, Springer, 2013, Theorem 1.8.10 and Corollary 1.8.13,
as cited in the source.
\bibitem{BDL} F. Buekenhout, M. Dehon and D. Leemans,
\emph{An atlas of residually weakly primitive geometries for small groups},
Acad. Roy. Belg., 1999; retained examples as cited in the source.
\end{thebibliography}
